\subsection{A coherent global norm family after an explicit radicand change}

\begin{statement}{A.149}{Proved}
{One compatible completed family in the changed tower at eleven}
Put
\[
 \begin{gathered}
 K=\Q(\alpha),\quad \alpha^3=3,\quad U=\alpha^2-2,
 \quad \gamma=2+\alpha,\\
 a'=3\cdot11^{24},\qquad z=\alpha11^8,\qquad
 \beta=(1-\alpha)U\,\frac{11}{(2+\alpha)^3}.
 \end{gathered}
\tag{A.149.1}
\]
Choose compatible positive roots
\(\rho_n^{3^n}=a'\), \(\rho_{n+1}^3=\rho_n\).
Thus \(\rho_1=z\), and
\[
 K'_n=\Q(\rho_n),\qquad K'_1=K,
 \qquad \alpha=\rho_1/11^8.
\tag{A.149.2}
\]
Let \(S_n\) consist of the places of \(K'_n\) above
\(2,3,11,\infty\).
There exists a compatible family in this fixed tower
\[
 \begin{gathered}
 b_n\in\widehat{\mathcal O_{K'_n,S_n}^{\times}}_{(3)}
                  \quad(n\geq1),\qquad b_1=\beta,\\
 N_{K'_{n+1}/K'_n}b_{n+1}=b_n,
                    \qquad N_{K'_n/\Q}b_n=-2.
 \end{gathered}
\tag{A.149.3}
\]
All transition maps are the actual field norms on the
displayed pro-three completed unit groups.
In particular, the three-primary one-variable obstruction
\(o_{\rm pol}(a',-2)\) of A.77--A.81 is zero.
The proof gives its coefficient and Shapiro maps.

The changed coefficient satisfies \(a'/3=11^{24}\).
The exponent-twenty-four old cover has both its rational
point of centre three and a unit point over \(\Q_{11}\).
These are finite-coordinate assertions.  The theorem
does not assert that \((a',-2)\) is an ordinary completed
unit-equation pair, or that it lifts to the full metabelian
quotient.  It establishes compatibility in this specified
changed global norm tower, including the prescribed
first representative \(\beta\).
\end{statement}
\begin{proof}
\emph{A ray-class calculation and its inertia consequence.}
Set \(k=\Q(\mu_3)\).
Its class number is one, by the Minkowski computation in
A.138.  The prime eleven is inert in \(k\), since
\(11\equiv-1\pmod3\); write \(\mathfrak q=(11)\).
Its residue field has order 121, and
\[
 v_3(121-1)=1,
 \qquad
 \mu_3\longrightarrow(\mathcal O_k/\mathfrak q)^\times
        \text{ maps onto its three-primary subgroup}.
\tag{A.149.4}
\]
The ray-class exact sequence
\[
 \mathcal O_k^\times\longrightarrow
 (\mathcal O_k/\mathfrak q^e)^\times
 \longrightarrow\operatorname{Cl}_{\mathfrak q^e}(k)
 \longrightarrow\operatorname{Cl}(k)\longrightarrow0
\tag{A.149.5}
\]
therefore has ray class group of order prime to three,
for every \(e\geq1\).
Indeed, passage from \(\mathfrak q\) to
\(\mathfrak q^e\) adds only eleven-primary groups.
The exact sequence and its field interpretation are
\cite[V, Theorems 1.7, 3.5 and 3.6]{MilneCFT}.
Consequently there is no nontrivial abelian
three-extension of \(k\) ramified only at
\(\mathfrak q\).

Let \(F/k\) be any finite Galois three-extension
ramified only at the primes above three and eleven.
Let \(I\) be inertia at one prime over three and
let \(\langle I\rangle^{\rm normal}\) be its normal
closure in \(G=\operatorname{Gal}(F/k)\).
The quotient field is ramified only at eleven.
Any nontrivial finite three-group has a nontrivial
abelian quotient, contrary to (A.149.5), so that
\(\langle I\rangle^{\rm normal}=G\).
If \(I\) were proper, it would lie in a maximal
subgroup of the finite three-group \(G\), which is
normal of index three.  Its normal closure would
then also be proper.  Hence
\[
                          I=G.
\tag{A.149.6}
\]
Thus every such extension is totally ramified at its
unique three-prime.

Apply this to
\[
                 L_n=\Q(\rho_n,\mu_{3^n}),\qquad n\geq1.
\tag{A.149.7}
\]
These are finite Galois three-extensions of \(k\),
ramified only at three and eleven, because the
cyclotomic degrees over \(k\) and the Kummer degrees
are powers of three.
Their three-primary Hilbert class fields are again
finite Galois three-extensions of \(k\), with the
same possible ramification: the Hilbert extension
is preserved by every automorphism of \(L_n/k\).
Equation (A.149.6) makes the Hilbert extension
totally ramified at three, whereas it is unramified
over \(L_n\).  The local and global degrees force
it to be trivial.  We have proved
\[
 \operatorname{Cl}(L_n)\{3\}=0,
 \qquad D_3(L_n/K)=\operatorname{Gal}(L_n/K)
                  \quad(n\geq1).
\tag{A.149.8}
\]
The second identity follows because \(k/\Q\) is
also totally ramified at three, and \(K\subset L_n\).
The same Hilbert base-change argument as A.147 gives
trivial three-primary class groups for every subfield
of \(L_n\), including \(K'_n\).

\emph{The specified target and its local norm families.}
The exact cubic norm gives
\[
 N_{K/\Q}(2+\alpha)=11,
 \quad N_{K/\Q}(1-\alpha)=-2,
 \quad N_{K/\Q}U=1,
 \quad N_{K/\Q}\beta=-2.
\tag{A.149.9}
\]
Let \(P\) and \(Q\) be respectively the degree-one
and degree-two primes of \(K\) above eleven.
The polynomial \(T^3-3\) has linear root 9 modulo
eleven and quadratic factor
\(T^2-2T+4\).
Since \(\gamma\) is integral of norm eleven and
vanishes on the linear factor,
\(\operatorname{div}(\gamma)=P\).
If \(P_2\) is the degree-one prime above two,
then
\[
                  \operatorname{div}(\beta)=P_2-2P+Q.
\tag{A.149.10}
\]
The coefficient two in \(NQ=2(11)\) makes the
eleven-part have norm zero, as required by
(A.149.9).

Put \(E_j=K(\rho_{j+1})=K'_{j+1}\) for \(j\geq0\),
so \(E_0=K\) and \(\rho_{j+1}^{3^j}=z\).
At three, the normalized valuation of \(z\) in
\(\Q_3(\alpha)\) is one.
Thus \([E_j:K]=3^j\), and A.140 supplies a
compatible actual local unit norm family for
the unit target \(\beta\).

At \(P\), the base field is \(\Q_{11}\), the
radicand valuation is eight, and
\(v_3(11-1)=0\).
The exact tame norm computation A.131 therefore
makes every element a norm at every depth,
including the valuation-minus-two target here.
Its compact fixed-valuation fibres give an actual
compatible local family.

At \(Q\), the base field is the unramified
quadratic extension of \(\Q_{11}\), with residue
field
\[
 \F_{121}=\F_{11}[t]/(t^2-2t+4).
\tag{A.149.11}
\]
Let \(\zeta=t^{40}\); it has order three.
For the cubic residue character defined by
\(w^{40}=\zeta^{\lambda(w)}\), exact arithmetic
in this field gives
\[
 \lambda(t)=1,\qquad \lambda(1-t)=0,
                       \qquad\lambda(t^2-2)=2.
\tag{A.149.12}
\]
The target has valuation one, and its unit
character is two, since \(\gamma^3\) contributes
zero character.
The radicand has valuation eight and unit
character one.  Its full stable symbol is
\[
                           8\cdot2-1=0\pmod3.
\tag{A.149.13}
\]
Here is the local norm computation over this
quadratic base, retaining its residue field.
At degree \(3^j\) choose integers \(A,B\) with
\(8A+3^jB=1\).
The uniformizer \(\rho_{j+1}^{A}11^B\) has norm
\(11\alpha^A\).
Unit norms have residues precisely
\((\F_{121}^{\times})^{3^j}\); all principal
units are scalar norms because their
\(3^j\)-power map is invertible.
Since \(v_3(120)=1\), the norm condition at
every positive depth is exactly (A.149.13).
Every fibre has fixed normalized valuation,
is compact, and is nonempty at each depth.
Taking the norm of an element at a largest
finite height solves every finite list of
compatibility equations.  Compactness proves
an actual compatible family at \(Q\).

At two, the local algebra of \(K\) has its
degree-one \(\Q_2\) factor and an unramified
quadratic factor.  On the first, the unit
radicand has compatible roots of every
three-power order in \(\Q_2\).
Put the target on the resulting degree-one
factor of every \(E_j\)-algebra and one on
the other factors.
On the quadratic factor the target is a unit
and all root algebras are unramified, so
their unit norm maps are onto.
At every finite place away from \(2,3,11\),
both target and radicand are units and the
same unramified-unit argument applies.
At real places use the odd polynomial's
compatible real branch, and at complex
places use a compatible complex branch.
We have thus constructed compatible local
norm families at every place, on the
actual etale algebras \(E_j\otimes_KK_v\).

\emph{Finite-coefficient duality with the full three-decomposition map.}
Let \(S\) be the finite set of places of \(K\)
above \(2,3,11,\infty\), and let \(G_{K,S}\)
be the Galois group unramified outside \(S\).
Consider a finite three-primary \(G_{K,S}\)-module
\(A\) whose action is trivial over some \(L_m\).
For a class
\[
 [c]\in\ker\!\left(H^1(G_{K,S},A)
               \longrightarrow\prod_{v\in S}H^1(K_v,A)\right),
\tag{A.149.14}
\]
choose a continuous cocycle and local witnesses
\(a_v\in A\) with \(c(g)=g a_v-a_v\) on
each selected local decomposition group.
Over \(L_m\), the cocycle is an abelian
three-primary character.
It is unramified outside the places above \(S\),
and its restrictions at those places are zero:
the group over \(L_m\) fixes every \(a_v\).
The same assertion holds at their conjugates
by the cocycle identity.
It therefore defines an everywhere unramified
abelian three-extension of \(L_m\), so is zero
by (A.149.8).
The cocycle now factors through
\(\operatorname{Gal}(L_m/K)\).
Its full three-decomposition group is that
entire quotient; the witness \(a_3\) thus
satisfies \(c(g)=g a_3-a_3\) for every
\(g\in G_{K,S}\).
This proves that (A.149.14) is zero.

If \(M\) is a finite three-primary module
split by the union \(L_\infty=\bigcup_mL_m\),
then its finite Cartier dual
\(M^D=\operatorname{Hom}(M,\mu_{3^\infty})\)
is also split by some \(L_m\).
Finite Poitou--Tate duality
\cite[I, Theorem 4.10(a)]{BIR_MilneADT}
identifies the two finite localization kernels
as perfect duals.  The preceding calculation
therefore proves the injective actual map
\[
 H^2(G_{K,S},M)\lhook\joinrel\longrightarrow
                         \prod_{v\in S}H^2(K_v,M).
\tag{A.149.15}
\]
This is a finite-set cohomology assertion
with the specified splitting fields.

\emph{The complete coefficient module and removal of the finite-set limits.}
Write \(X=1+T\), \(\Lambda=\Z_3[[T]]\), and let
\(a:G_{K,S}\to\Z_3(1)\) be the Kummer cocycle of
the fixed \(S\)-unit \(z\).
The actual coefficient sequence of A.81 over \(K\) is
\[
 0\longrightarrow D_z\xrightarrow{\,T\,}V_z
       \xrightarrow{\epsilon}\Z_3(1)\longrightarrow0,
 \qquad D_z=V_z=\Lambda\quad\text{as underlying modules},
\tag{A.149.16}
\]
with actions
\[
 \begin{aligned}
 g\cdot v&=\chi(g)X^{a(g)}v(X^{\chi(g)}-1),\\
 g\cdot d&=\chi(g)X^{a(g)}
     \frac{X^{\chi(g)}-1}{X-1}
                         d(X^{\chi(g)}-1).
 \end{aligned}
\tag{A.149.17}
\]
The notation in (A.149.17) regards \(v,d\) as
power series in \(T\).
Multiplication by \(T\) is equivariant, and
the sequence splits continuously as underlying
\(\Z_3\)-modules.
All its finite coefficient actions factor
through \(L_\infty/K\), since that field contains
all roots of \(z\) and all three-power roots
of unity.
In particular the stable finite modules
\[
                  D_{r}=D_z/(3^r,T^r)D_z
\tag{A.149.18}
\]
are split by suitable actual \(L_m\).

Let \(o=\delta[\beta]\in H^2_{\rm cts}(G_{K,S},D_z)\)
be the connecting class of (A.149.16).
Fix \(r\), and choose \(j\) sufficiently large
that \((X^{3^j}-1)/(X-1)\) lies in
\((3^r,T^r)\); the binomial bound in A.81
provides this exact inclusion.
The permutation quotient (A.149.20) and its
augmentation kernel (A.149.22) give a morphism
of the extension (A.149.16) to the extension
with kernel \(D_r\).
At each place \(v\in S\), one norm preimage
at depth \(j\) supplies a local degree-one
class in that permutation module with
augmentation \([\beta]\).
Its connecting class, and therefore the
image of \(o\) in \(H^2(K_v,D_r)\), is zero.
This finite-depth comparison already suffices;
the stronger compatible local families were
proved above.
By (A.149.15), the image of \(o\) in
\(H^2(G_{K,S},D_r)\) is zero for every \(r\).

The finite groups \(H^1(G_{K,S},D_r)\) are finite
by \cite[I, Corollary 4.15]{BIR_MilneADT}.
Their inverse system is Mittag--Leffler.
The coefficient maps \(D_{r+1}\to D_r\) are onto,
and their set sections lift continuous cochains.
The continuous Milnor sequence therefore gives
\[
 H^2_{\rm cts}(G_{K,S},D_z)
       \xrightarrow{\ \sim\ }
                 \varprojlim_rH^2(G_{K,S},D_r).
\tag{A.149.19}
\]
Consequently \(o=0\).
This explicitly removes the coefficient and
augmentation limits over the fixed finite set
\(S\); it does not assume their removal over an
unrestricted absolute Galois group.
One can equally first take the finite free
augmentation quotients \(D_z/T^r\): their
finite-coefficient \(H^1\) groups give compact
integral \(H^1\), whose continuous inverse
system has zero algebraic \(\varprojlim^1\)
by the compact finite-intersection argument.
Both procedures give the same projections in
(A.149.19).

\emph{Shapiro gives the same global norm family at every level.}
The vanishing of \(o\) supplies
\(h\in H^1(G_{K,S},V_z)\) mapping to \([\beta]\).
For \(j\geq0\), the actual permutation quotient is
\[
 V_{z,j}=V_z/(X^{3^j}-1)V_z
          =\Z_3(1)[\operatorname{Hom}_K(E_j,\overline K)].
\tag{A.149.20}
\]
Indeed, labeling the roots by
\(\zeta_{3^j}^i\rho_{j+1}\), the Galois action on
the index is \(i\mapsto\chi(g)i+a(g)\), exactly
the action in (A.149.17) on the basis \(X^i\).
The topology also retains the exact comparison
\(V_z=\varprojlim_j V_{z,j}\).
Indeed, the ideals \((X^{3^j}-1)\) are nested
and tend to zero in the \((3,T)\)-adic topology
by the same binomial bound.  Their intersection
is zero, and compactness of \(\Lambda\) makes
every compatible sequence of quotient cosets
have a common element.  Each quotient is finite
free over \(\Z_3\); its onto transition map
has a continuous underlying \(\Z_3\)-linear
section.  Hence the continuous cochain systems
used here are degreewise onto.
Continuous Shapiro and Kummer theory identify
\[
 H^1(G_{K,S},V_{z,j})
  =H^1(G_{E_j,S_{j+1}},\Z_3(1))
  =\widehat{\mathcal O_{E_j,S_{j+1}}^{\times}}_{(3)}.
\tag{A.149.21}
\]
For the last equality, the finite Kummer
sequence has the \(S_{j+1}\)-class group torsion
as its cokernel; its inverse Tate module is
zero because the class group is finite.
Here its three-primary part also vanishes
by (A.149.8) and Hilbert base change.
The coefficient map
\(V_{z,j+1}\to V_{z,j}\) sends each root to its
cube.  Under Shapiro its corestriction is
precisely \(N_{E_{j+1}/E_j}\); augmentation is
\(N_{E_j/K}\).
The projections of the single class \(h\)
therefore give (A.149.3), beginning with
\(b_1=\beta\), and their norm to \(\Q\) is
\(-2\) by (A.149.9).

For the comparison with augmentation coordinates,
multiplication by \(T\) identifies the kernel
of the augmentation in (A.149.20) with
\[
 D_z\Big/\left(\frac{X^{3^j}-1}{X-1}\right)D_z.
\tag{A.149.22}
\]
The binomial bound of A.81 gives a coefficient
map from this kernel to every \(D_r\) once
\(j\) is large enough.  These are the actual
maps linking the permutation norms and the
finite augmentation tests used above.
Finally include level \(\Q\), with target
\(-2\), in (A.149.3).  The root-torsor Shapiro
map for \(a'\) over \(\Q\) then gives a lift
of \([-2]\) to its module \(V_{a'}\): the
degree-zero permutation cohomology groups
are zero, so the inverse family of degree-one
classes is the image of one continuous class.
Its connecting class is precisely
\(o_{\rm pol}(a',-2)=0\).

For the retained old-cover test, \(2^{24}\) and
\(9^{24}\) are both \(-1/2=5\) modulo eleven.
Hensel lifting gives a unit solution with
\(X=1\) to
\(a'X^{24}-2Y^{24}=1\) over \(\Q_{11}\).
The rational solution \((X,Y)=(11^{-1},1)\)
has centre three.
These computations complete the asserted
finite-coordinate comparison and the
specified completed global compatibility.
\end{proof}

The exact certificate
\path{code/coherent_inert_radicand_update_certificate.py}
and its retained output verify the ray residue group,
all finite-field norm symbols, the old-cover roots,
the divisor norm weights, and the coefficient
projection bounds.  The class-field, duality and
inverse-limit steps are proved above, rather than
inferred from the finite certificate.
